\documentclass[11pt]{article}
\usepackage[margin=1in]{geometry}
\usepackage{amsmath,amssymb,amsthm,mathtools}
\usepackage{array,booktabs,microtype}
\usepackage{xcolor}
\usepackage[colorlinks=true,linkcolor=blue!45!black,citecolor=blue!45!black,urlcolor=blue!45!black]{hyperref}
\newtheorem{theorem}{Theorem}
\newtheorem{proposition}[theorem]{Proposition}
\newtheorem{lemma}[theorem]{Lemma}
\newtheorem{corollary}[theorem]{Corollary}
\DeclareMathOperator{\CT}{CT}
\DeclareMathOperator{\End}{End}
\DeclareMathOperator{\Tr}{Tr}

\title{The primitive balanced-degree spectrum\\in the two-pair Image problem}
\author{Roy van Rijn}
\date{}

\begin{document}
\maketitle

\begin{abstract}
Over every characteristic-zero field, the balanced degrees admitting radially
primitive, absolutely nonpower two-pair Image counterexamples are precisely
\(4,6,7,8,\ldots\). All pure balanced phases through degree three are safe,
as are all pure radially primitive balanced quintics. One explicit integer
family realizes the failure degrees with a common marker of balanced degree one that detects
every positive power. We also prove that an abelian variety with
Rosati-compatible rank-two multiplication by a CM field has ordinary
primes of density one after finite extension if at most two conjugate
embedding pairs have Hodge multiplicities \((1,1)\). The full geometric
endomorphism algebra may be larger than the acting field. In the quintic
proof, this theorem supplies ordinary reduction for factors of cyclic
covers whose Frobenius coefficients detect nonzero polynomial moments.
\end{abstract}

\section{Introduction}

Let \(K\) be a field of characteristic zero. For
\(\xi=(\xi_1,\xi_2)\) and \(z=(z_1,z_2)\), define the contraction on monomials by
\(\mathcal E(\xi^\alpha z^\beta)=\partial_z^\alpha z^\beta\) and extend
linearly. The Image space is the sum of the images of
\(\xi_i-\partial_{z_i}\), \(i=1,2\), and equals \(\ker\mathcal E\)
\cite[Theorem~3.3]{image}.
A phase \(F\) is pure if \(\mathcal E(F^m)=0\) for all \(m\geq1\).
It is safe if for every fixed polynomial \(Q\), \(\mathcal E(QF^m)=0\)
for all sufficiently large \(m\). Balanced degree \(d\) means bidegree
\((d,d)\) in \((\xi,z)\). A nonzero balanced phase is radially primitive
if \(R=\xi_1z_1+\xi_2z_2\) does not divide it. It is absolutely nonpower if
it is no nonzero scalar multiple of a proper power over an algebraic closure.

\begin{theorem}\label{thm:spectrum}
Let \(K\) have characteristic zero.
\begin{enumerate}
\item Every pure two-pair balanced phase of degree at most three is safe.
\item Every pure, radially primitive two-pair balanced phase of degree five
      is safe.
\item In degree four and every degree at least six there are explicit integer,
      radially primitive, absolutely nonpower two-pair balanced phases \(F\)
      with \(\mathcal E(F^m)=0\) and \(\mathcal E(ZF^m)\ne0\) for every
      \(m\geq1\), with the common linear marker \(Z=\xi_1z_2\) in the
      balanced coordinates.
\end{enumerate}
Consequently the primitive, absolutely nonpower balanced counterexample
degrees are \(4,6,7,8,\ldots\).
\end{theorem}

The Special Image Conjecture asserts safety for every pure phase in the
original polynomial ring \cite[Conjecture~3.1]{image}. A single quartic
witness already refutes that assertion in two pairs. The question here is
which balanced degrees can carry a failure with the radial factor
removed. Allowing radial factors, the balanced failure spectrum starts
at four and contains every higher degree. Indeed, multiplying the quartic by
\(R^{d-4}\) leaves its restriction to \(R=1\) unchanged, and
\eqref{eq:spectrum-contraction} changes pure and marked contractions only
by nonzero factorial ratios. Among phases divisible by \(R\), the failure
degrees start at five, since a divisible quartic restricts to sphere
degree at most three. Radial primitivity is the condition that degree in the
original balanced ring agrees with least degree on the sphere.

Long's Hopf identity gives the quartic seed \cite{longcompact}. We give a
joint endpoint identity that lets the seed be multiplied by selected
quadratic and cubic factors while retaining the same marker. This
produces degrees \(4+2a+3b\) for \(a,b\geq0\). Proving that five is genuinely
missing requires controlling arbitrary primitive quintic coefficients.

An abelian variety of dimension \(g\) has ordinary reduction at a good
prime when its Newton slopes consist of \(g\) zeros and \(g\) ones.

\begin{theorem}[Ordinary reduction with rank-two CM multiplication]\label{lem:rank-two}
Let \(J\) be an abelian variety of dimension \(2e>0\) over a number field \(L\).
Suppose a CM field \(E\) of degree \(2e\) embeds in
\(\End^0(J_{\overline L})\) and a polarization induces complex conjugation
on \(E\) by its Rosati involution. For each conjugate pair of embeddings of
\(E\), consider the multiplicities of the two embeddings on holomorphic
one-forms. Suppose at most two pairs have multiplicities \((1,1)\). Then
there is a finite extension of the number field over which the ordinary
primes of \(J\) have natural density one.
\end{theorem}

Theorem~\ref{lem:rank-two} concerns all abelian varieties with the stated
action and signature. Its relative-rank-two range includes signatures with
two balanced pairs. At a prime split in \(E\), a nonordinary balanced
pair forces divisibility of the Frobenius trace at \emph{both} places of
that pair. With at most two balanced pairs this supplies the remaining
divisibility beyond the contribution from the definite pairs. The Weil
bound puts the normalized trace norm in a finite set of integers. After
adjoining enough torsion, the norm is congruent to its largest possible
value, forcing equality. Equality contradicts the two local divisibilities
at the nonordinary balanced pair.

The link to polynomial moments comes from the interval integration after
angular averaging. Integration divides a term of height \(j\) by \(j+1\).
At special moment orders and primes, the largest height is \(p-1\), so
one prime denominator isolates one coefficient of a power of the leading
face. The same coefficient occurs as the Frobenius scalar on a
one-dimensional eigenspace of the coherent cohomology of a cyclic cover.
An ordinary Jacobian makes that scalar nonzero. This rules out purity even
though individual terms of the original moment can cancel in
characteristic zero.

The proof of safety starts with the leading binary form on the conic at
infinity. If its intersection orders are all small, the volume extends
across a proper resolution and a nonzero constant exponential period is
impossible. The borderline intersection order gives a single nonzero
coefficient under the prime test. Higher order leaves finitely many
supporting faces, where the cyclic-cover argument applies except in
explicit cases that reduce to safe phases by a change of coordinates.
The quadratic-height analysis below uses the general root and face tests
from Section~\ref{sec:quintic}. Those tests are independent of the collision
analysis, which uses quadratic height to finish the remaining cases.

\section{Reduction to spherical moments}\label{sec:sphere}

Contraction kills the image of each \(\xi_i-\partial_{z_i}\). Applying
this operator to \(\xi^{\alpha-e_i}z^\beta\) replaces one outer factor
\(\xi_i\) by \(\partial_{z_i}\) modulo its image. Induction on the
\(\xi\)-degree shows that every \(H-\mathcal E(H)\) belongs to the sum
of the two images, proving the kernel identity.

Introduce the four balanced bilinears
\[
 R=\xi_1z_1+\xi_2z_2,\quad Z=\xi_1z_2,\quad
 W=2\xi_2z_1,\quad T=\xi_1z_1-\xi_2z_2,
 \qquad T^2+2ZW=R^2.
\]
They generate the balanced subring, with the displayed relation.
Generation follows by pairing one \(\xi_i\) with one \(z_j\). The normal
monomials modulo the relation in homogeneous degree \(d\) number
\((d+1)^2\), the dimension of the original bidegree-\((d,d)\) space,
so there is no further relation. Restrict at \(R=1\), writing
\((x,y,t)=(Z,W,T)\) there. The sphere algebra and its trace are
\[
 A_K=K[x,y,t]/(t^2+2xy-1),\qquad
 \Lambda(h)=\operatorname{CT}_x\frac12\int_{-1}^1
    h\biggl(x,\frac{1-t^2}{2x},t\biggr)\,dt.
\]
The integral is the rational polynomial rule with value \(1/(j+1)\)
on \(t^j\) for even \(j\), and zero for odd \(j\), including the factor
\(1/2\). If \(H\) has bidegree \((d,d)\) and sphere restriction \(h\), then
\begin{equation}\label{eq:spectrum-contraction}
 \mathcal E(H)=(d+1)!\,\Lambda(h).
\end{equation}
To check this, take an original monomial with \(\xi_1\)-exponent \(i\)
and \(z_1\)-exponent \(j\). Both sides vanish for \(i\ne j\) by angular
weight. For \(i=j\) the sphere integral is the beta integral
\(i!(d-i)!/(d+1)!\), and its contraction is \(i!(d-i)!\).

The least sphere degree of \(h\) is \(d\) exactly when \(R\) does not
divide \(H\). Indeed homogenizing the normal monomials \(x^at^b\),
\(y^ct^b\) with \(c>0\) recovers \(H\), and each has filtration degree
\(a+b\) or \(c+b\), respectively. We use \emph{pure} for the condition
\(\Lambda(h^m)=0\) for every \(m\geq 1\) and \emph{safe} for the assertion
that \(\Lambda(gh^m)=0\) eventually for each fixed \(g\in A_K\).

Sphere safety gives the required vanishing against arbitrary original
polynomial markers. For a bihomogeneous marker \(Q\) of bidegree
\((a,b)\), contraction against \(H^m\) is zero if \(a>b\). For \(b\ge a\),
its output has degree \(b-a\) in \(z\), and for every \(|\gamma|=b-a\),
\[
 \partial_z^\gamma\mathcal E(QH^m)
   =\mathcal E(\xi^\gamma QH^m).
\]
The marker on the right is balanced. There are finitely many such
derivatives, so their eventual vanishing by
\eqref{eq:spectrum-contraction} forces eventual vanishing of the
output. A general marker is the sum of finitely many bihomogeneous
ones. Conversely any fixed sphere marker with infinitely many nonzero
mixed moments lifts to a balanced marker by homogenization.

\section{Quadratic height and the degrees below four}\label{sec:quadratic-height}

We first treat phases whose Laurent expression after eliminating \(y\)
has degree at most two in \(t\). A coefficient of \(x^{-j}\), \(j>0\),
is divisible by \((1-t^2)^j\), so only a simple pole at \(x=0\) is possible.
With no pole, the angular core must vanish for purity by the interval
polynomial theorem \cite[Theorem~2.1]{muger-tuset}. All nonzero angular exponents are then positive,
giving safety. With a pole, scale the phase and write
\[
 h=x^{-1}(A-Bt^2+Ct),\qquad A,B,C\in K[x],\quad A(0)=B(0)=1,\quad C(0)=0.
\]
Put \(D=A+C^2/(4B)\) and \(M=\sqrt{D/B}\), with \(M(0)=1\).
In completing the square, set \(t=C/(2B)+Mz\). The endpoints in \(z\)
are congruent to \(\pm1\) modulo \(x\). The integrand \((1-z^2)^m\)
has contact order \(m\) there, so returning the endpoints to \(\pm1\)
changes its integral by a multiple of \(x^{m+1}\). Therefore
\[
 \Lambda(h^m)=\frac12\int_{-1}^1(1-z^2)^m\,dz\;[x^m]D^mM.
\]
The factor given by the integral is a nonzero rational number.
The conditions \([x^m]D^mM=0\) determine each successive coefficient of
\(M\) since \(D(0)=1\). Differentiating \(D^m\) shows that the unique
solution is \(M=1-xD'/D\). In particular
\begin{equation}\label{eq:quintic-quadratic-height}
 B=\frac{D^3}{(D-xD')^2}.
\end{equation}

Suppose the phase has least sphere degree at most five, so \(\deg B\le4\).
At any finite zero or pole \(\alpha\) of \(D\) of order \(r\ne0\),
we have \(\alpha\ne0\), \(\operatorname{ord}_\alpha(D-xD')=r-1\), and
\(\operatorname{ord}_\alpha B=r+2\). There are no other finite zeros of
\(B\) by \eqref{eq:quintic-quadratic-height}. If \(n\) is the degree of
\(D\) at infinity and \(s\) counts its distinct finite zeros and poles
over an algebraic closure, summing valuations gives \(\deg B=n+2s\).
When \(n\ne1\), the leading coefficient of \(D-xD'\) is \(1-n\)
times that of \(D\), hence \(\deg B=n\) and \(D\) is constant.
When \(n=1\), the bound \(\deg B\le4\) forces \(s=1\), so \(D\) is
linear and \(B=D^3\). Then \(B\) divides \(C^2\), and \(C=xD^2P\)
for a polynomial \(P\). Since \(A=D-x^2DP^2/4\) has degree at most six,
\(P\) has degree at most one. The phase has exact degree four, because
its \(x^{-1}Bt^2\) term has degree four and its other terms have degree
at most four.

In the constant case \(D=B=1\). With \(c=C/x\), the phase is
\(2y+ct-xc^2/4\), the image of \(2y\) under the polynomial shear
\((x,t,y)\mapsto(x,t+xa,y-at-xa^2/2)\) with \(a=-c/2\).
The derivation \(xa\partial_t\) in the literal Laurent chart kills the
trace by integration at \(t=\pm1\). Each boundary value is a polynomial
in \(x\) since \(y=0\) there, and the extra factor \(x\) kills its constant
term. The shear is its finite exponential on each regular polynomial,
so both directions preserve the trace and safety. Thus in quadratic
height every pure phase of least degree at most five is safe or has
degree four.

Here is the exclusion in degrees \(d=1,2,3\), using the leading-form
test proved in Section~\ref{sec:leading-threshold}. In a pure phase the
binary leading form has a root of multiplicity greater than \(d\).
Move the root to the point \(v=0\) in the parametrization
\((x,y,t)=(u^2,-v^2/2,uv)\). Every degree-\(d\) normal monomial then
has negative angle \(i\). If the phase has no positive angle it is
safe by the interval theorem and one-sided support. Otherwise, use the
supporting-line and two-term tests of Section~\ref{sec:face-support} with
\(d\) in place of five. The allowed nodes have \(-d\le i\le d-1\) and
\(\max(0,-2i)\le j\), with \(i+j\le d\) for \(i<0\) and
\(i+j\le d-1\) otherwise. A crossing line \(j+\rho i=H\) has
\(\rho>1\) and \(H<d\).

Joining opposite-sign nodes in that range gives no line with three
or more nodes for \(d=1,2\). For \(d=3\) the only such lines are
\((\rho,H)=(2,2)\), with angles \(-1,0,1\), and
\((\rho,H)=(3/2,3/2)\), with angles \(-3,-1,1\). The former has
quadratic height after interchanging \(x,y\), so is safe by the
preceding argument. The latter has face \(Bx+Ayt+Cy^3\) with all
three coefficients nonzero. Its order-three cover has branch
exponents \(2,2,2\) when the face quadratic has distinct roots, so its
elliptic Jacobian supplies the ordinary-prime obstruction of
Section~\ref{sec:cyclic-face}. If the quadratic repeats, then
\(A^2+2BC=0\). The face is \(B\) times the leading part of
\(x\mapsto x-cyt-c^2y^3/2\) with \(c=-A/B\), in the shear fixing
\(y\) and sending \(t\) to \(t+cy^2\). Its inverse preserves the weights
\((3,1,2)\) on \((x,y,t)\) and removes the cubic part, leaving
degree at most two. This last phase is safe by the preceding
low-degree cases. Degree zero is immediate since a nonzero constant
has nonzero first moment. The sharper mixed-moment bound for cubics
is proved in \cite{cubicsphere}.

\section{Explicit primitive witnesses}\label{sec:witnesses}

The following construction uses the quartic endpoint mechanism of
Long \cite[Theorem~3.1]{longcompact}. A joint endpoint identity allows
an additional cubic factor while keeping all positive power moments
zero. Set
\[
 U=1+x,\qquad w=tU,\qquad
 e=\frac{U}{x}(1-w^2)=U\bigl(2y-(2+x)t^2\bigr),\qquad
 g=e+w^2=U(2y-t^2).
\]
For integers \(r\ge1\) and \(s\ge0\) put
\[
 J_{r,s}(V)=\frac12\int_{-V}^V z^s(1-z^2)^r\,dz.
\]
Here the endpoints are evaluated in a polynomial primitive.
Both endpoint integrands vanish to order at least \(r\) at \(V=1\),
so \(J_{r,s}(V)-J_{r,s}(1)\) is divisible by \((V-1)^{r+1}\).
The substitution \(z=tU\) gives
\[
 \Lambda(e^r w^s)=[x^r]U^{r-1}J_{r,s}(U)=0,\qquad
 \Lambda(xe^r w^s)=[x^{r-1}]U^{r-1}J_{r,s}(U)=J_{r,s}(1).
\]
In each coefficient the endpoint error is of order at least \(r+1\),
and the remaining factor \(U^{r-1}\) is monic of degree \(r-1\).
Thus in all polynomials divisible by \(e\) in \(K[e,w]\), the trace
vanishes, and the \(x\)-marked trace agrees with the interval
substitution \(e=1-z^2\), \(w=z\). In particular, for \(a,b\ge0\),
the phase \(e(1+w)^a g^b\) is pure. Its \(x\)-marked \(m\)th power
has trace
\begin{equation}\label{eq:spectrum-beta}
 \frac12\int_{-1}^1(1-z^2)^m(1+z)^{am}\,dz
 =2^{(a+2)m}\frac{((a+1)m)!\,m!}{((a+2)m+1)!}.
\end{equation}
The factor \(g^b\) disappears from this integral since \(g=e+w^2=1\)
under the marked substitution. Equation \eqref{eq:spectrum-beta}
follows by putting \(z=2u-1\) and integrating the beta monomial
\(u^{(a+1)m}(1-u)^m\).

The corresponding integer forms can be written using
\[
 L=R+Z,\quad D=2WR^2-(2R+Z)T^2,\quad E=2WR-T^2,\quad
 A_2=R^2+TL,\quad G_3=LE.
\]
For \(a,b\ge0\) set \(F_{a,b}=LD\,A_2^aG_3^b\) and \(d=4+2a+3b\).
Each factor is balanced homogeneous, and the sphere restriction is
\(e(1+w)^a g^b\). Equations \eqref{eq:spectrum-contraction} and
\eqref{eq:spectrum-beta} give
\[
 \mathcal E(F_{a,b}^m)=0,\qquad
 \mathcal E(ZF_{a,b}^m)=(dm+2)!\,2^{(a+2)m}
             \frac{((a+1)m)!\,m!}{((a+2)m+1)!}\ne0
 \qquad(m\ge1).
\]
At \((\xi_1,\xi_2,z_1,z_2)=(1,1,1,-1)\) we have \(R=0\) and
\(F_{a,b}=-4(-2)^a4^b\ne0\), which proves radial primitivity.
Specializing \(\xi_1=z_1=1\), \(z_2=0\), \(\xi_2=u\) yields
\(2^a(4u-2)(4u-1)^b\), with a simple zero at \(u=1/2\). Any scalar
multiple of a proper power would specialize to zero or to a scalar
multiple of a proper power. Hence these forms remain nonpowers over
every field extension. Taking \(a=b=0\) realizes degree four,
taking \(b=0\) realizes every even degree at least six, and taking
\(b=1\) realizes every odd degree at least seven.

\section{Safety in least sphere degree five}\label{sec:quintic}

Write a sphere phase in the normal monomials \(x^i t^j\) with \(i\ge0\)
and \(y^i t^j\) with \(i>0\). Its leading symbol of degree \(d\), restricted
to the cone \(t^2+2xy=0\), gives a binary form \(h_d(u,v)\) of degree \(2d\)
by \((x,y,t)=(u^2,-v^2/2,uv)\). This identification sends the degree-\(d\)
normal monomials to distinct binary monomials. The \(\mathrm{SL}_2\) action
on \((u,v)\) induces linear isometries of the sphere, by congruence on
\(\left(\begin{smallmatrix}x&t\\t&-2y\end{smallmatrix}\right)\).
They preserve \(\Lambda\), since the sphere trace in each even homogeneous
degree is proportional to the Gaussian trace with covariance inverse to
\(t^2+2xy\), and the odd-degree traces are zero. This proportionality
follows by integrating each normal monomial. In particular we may move a
geometric root of \(h_d\) to either coordinate point after scalar extension.

\subsection{The leading root threshold}\label{sec:leading-threshold}

Suppose every root of \(h_d\) has multiplicity less than \(d\). Embed the
finitely generated coefficient field in \(\mathbb C\). Compactify the sphere
by the smooth projective quadric with boundary conic \(D\). Its affine
residue volume \(\omega\) has divisor \(-2D\) by adjunction. In local
coordinates \((u,z)\) with \(D=\{u=0\}\), write the phase as \(F/u^d\).
At a boundary root of multiplicity \(r>0\), Weierstrass preparation writes
\(F\) as a unit times a monic degree-\(r\) polynomial in \(z\). For fixed
\(u\), the set \(|F|\le\delta\) lies in the union of \(r\) planar disks of
radius \(O(\delta^{1/r})\). Thus on \(|f|\le M\) the squared volume density
has local integral bounded by a constant times
\[
 \int_{|u|<\epsilon}|u|^{-4+2d/r}\,d^2u,
\]
which is finite since \(r<d\). Away from the roots, the region \(|f|\le M\)
stays away from \(D\).

Resolve the pencil and remove its fibre over infinity. On the resulting
smooth surface \(Y\), \(f:Y\to\mathbb A^1\) is proper. The pulled-back
volume form is algebraic and holomorphic on \(Y\), since a divisorial pole
would give an infinite squared integral in a bounded-\(f\) chart.
Take the compact cycle \(\Gamma=\{y=\overline x,\ t\in\mathbb R,\
2|x|^2+t^2=1\}\), which lies in the original affine chart. For
\(-1<t<1\), write \(x=\sqrt{(1-t^2)/2}\,e^{i\theta}\). In these coordinates
the residue form is a nonzero constant multiple of \(d\theta\wedge dt\),
using its expression \(\pm dx\wedge dt/(2x)\) in the chart \(x\ne0\).
Consequently \(\int_\Gamma\omega\ne0\), and integrating \(f^m\omega\)
against \(\Gamma\) gives \((\int_\Gamma\omega)\Lambda(f^m)\).
Since the exponential series converges uniformly on the compact cycle on
bounded \(\tau\)-sets, purity would give the nonzero constant period
\(\int_\Gamma e^{-\tau f}\omega=\int_\Gamma\omega\).
Lemma~\ref{lem:proper-volume}, proved below, excludes this. Thus a pure
phase has a leading root of multiplicity at least \(d\).

Here is the prime-isolation rule used in the remaining cases. Work over
a number field at a good odd prime \(p\) where all coefficients are
integral. If the angular constant term of \(h^m\) is a polynomial
\(C_m(t)\) of degree at most \(p-1\), then
\begin{equation}\label{eq:prime-isolation}
                 p\Lambda(h^m)\equiv [t^{p-1}]C_m(t)\pmod{\mathfrak p}.
\end{equation}
Indeed the integral kills odd heights and divides even-height coefficients
by \(j+1\). The only denominator divisible by \(p\) in the given range is
the denominator at the even height \(p-1\).

An exact multiplicity \(d\) is also impossible under purity. Move that
root to \(u=0\). In the literal Laurent substitution \(y=(1-t^2)/(2x)\),
every monomial \(x^i t^j\) has \(i+j\le d\), and the part on equality is
\(t^d L(x/t)\) for \(L(z)=z^{-d}h_d(z,1)\). With the chosen root, \(L\) is
a polynomial with nonzero constant. At a good prime \(p=1\bmod d\) put
\(m=(p-1)/d\). The coefficient at height \(p-1\) is
\(\operatorname{CT}L^m=L(0)^m\), a unit, so \eqref{eq:prime-isolation}
applies. This argument may be made after
one number-field specialization preserving the leading identity and its
specified nonzero constants. Indeed a finitely generated coefficient
algebra over \(\mathbb Q\), localized to keep these constants nonzero, has a
maximal-ideal quotient that is a number field, and all rational moment
identities specialize. We use the same specialization argument in the
prime tests below, keeping the required distinct roots and support
coefficients nonzero. Primes in the progression \(1\) modulo a positive
integer are infinite, for example by the elementary cyclotomic argument.

\subsection{Supporting faces for multiplicity at least six}\label{sec:face-support}

Let \(d=5\) and move a root of multiplicity at least six to \(v=0\).
All degree-five normal terms now have negative angular exponent, since
their \(v\)-exponents are \(5-i\) for angular exponent \(i\). If the phase
has no positive angular terms, purity forces its zero-angle core to
vanish by the interval moment theorem \cite[Theorem 2.1]{muger-tuset}.
The negative support then gives eventual mixed vanishing for every fixed
marker. Otherwise take the upper supporting line \(j+\rho i=H\) of the
literal Laurent support over angular exponent zero. A convex combination
of nodes reaching angle zero has height less than five, since all
equality nodes of \(i+j\le5\) have negative angle. This gives \(H<5\).
For any degree-five node \(j+\rho i=5+(\rho-1)i\le H\) with \(i<0\),
so \(\rho>1\). In negative angle \(-b\), only the highest Laurent term
of each \(y^b t^r\) can lie on the upper boundary. Thus the possible face
nodes satisfy
\[
 -5\le i\le4,\qquad \max(0,-2i)\le j\le
 \begin{cases}5-i,&i<0,\\4-i,&i\ge0.\end{cases}
\]
A face with no opposite signs must contain a core node at angle zero.
Its height \(H\) is a positive integer, since the support contains nodes
of both signs and all negative-angle nodes have positive height.
At \(m=(p-1)/H\) for good primes \(p\equiv1\bmod H\), the core coefficient
to the \(m\)th power occupies the top angular-constant slot. Equation
\eqref{eq:prime-isolation} excludes it. For a two-node crossing face, write its
angles \(-b,a\) and heights \(j_-,j_+\). Set \(M=aj_-+bj_+\) and choose
a good prime \(p\equiv1\bmod M\), writing \(k=(p-1)/M\).
At exponent \(m=(a+b)k\), angular cancellation uses \(ak\) negative and
\(bk\) positive terms, giving height \(p-1\). The coefficient is the
product of \(\binom{(a+b)k}{ak}\) and a unit monomial in the two
coefficients. The factorials have arguments below \(p\) because
\(M\ge2ab\ge a+b\). Equation \eqref{eq:prime-isolation} excludes this case.

The remaining lines are listed below. To obtain the list, take the line
through every pair of allowed nodes of opposite signs, keep those with
\(\rho>1\), \(H<5\), at least one degree-five node on or below the line,
and at least three equality nodes. An actual crossing face occupies both
signs and at least three listed nodes.
\[
\begin{array}{c|c|l@{\qquad}c|c|l}
\rho&H&i&\rho&H&i\\ \hline
5/3&5/3&-5,-2,1&2&2&-3,-2,-1,0,1\\
3/2&5/2&-5,-3,-1,1&3/2&3&-4,-2,0,2\\
2&3&-2,-1,0,1&3&3&-1,0,1\\
4/3&10/3&-5,-2,1&3/2&7/2&-3,-1,1\\
4/3&11/3&-4,-1,2&5/4&15/4&-5,-1,3\\
4/3&4&-3,0,3&3/2&4&-2,0,2\\
2&4&-1,0,1,2&3/2&9/2&-1,1,3
\end{array}
\]

At \(\rho=H=2\), normal terms have weight at most two for weights
\((x,y,t)=(2,0,1)\). After exchanging \(x,y\), the literal Laurent height
is at most two. Regularity then allows at most a simple negative angular
pole. Section~\ref{sec:quadratic-height} proves safety for every such pure
phase of least sphere degree five.

\subsection{The cyclic-cover obstruction on a face}\label{sec:cyclic-face}

For a noninteger height \(H=n/q\) in lowest terms, the angular Laurent
polynomial of the face is \(z^{-b}P(z^q)\), where \(-b,a\) are its occupied
extreme angles. The table gives the same \(q\) as the denominator of \(\rho\).
In an integer-height row, first divide the angles by their common step if
the displayed angles are all multiples of two or three, and then put
\(q=1\) and \(n=H\). Angular compression leaves constant terms unchanged.
In these coordinates \(\deg P=(a+b)/q\), because both extreme terms
are occupied. For a prime \(p=nk+1\) and moment order \(qk\), the constant angular term has
height at most \(qkH=p-1\). A product contributing at that height uses
only equality nodes of the supporting line. The only interval coefficient
with denominator divisible by \(p\) is the one of height \(p-1\). Thus
multiplying the moment by \(p\) and reducing gives
\[
 [z^{bk}]P(z)^{qk}=[z^{p-1}]\bigl(z^{n-b}P(z)^q\bigr)^k.
\tag{*}\label{eq:face-scalar}
\]
The raw degree-\(2n\) branch form in the second expression has endpoint
exponents \(n-b,n-a\), and each root of \(P\) has its multiplicity multiplied
by \(q\). When all nonzero exponents are less than \(n\), the normalized
cyclic cover \(w^n=z^{n-b}P(z)^q\) detects \eqref{eq:face-scalar} in its
character-one \(H^1(O)\) line. In fact the character-one sheaf is \(O(-2)\)
and its \v{C}ech generator \(w/z\) maps to
\(wz^{-p}(z^{n-b}P(z)^q)^k\) under Frobenius. Projection to the line gives
exactly the selected coefficient. More generally the character-\(i\)
dimension is \(-1+\sum_j\{i a_j/n\}\) for branch exponents \(a_j\), by
dividing the local Kummer generator \(w^i\) by the branch coordinate to
power \(\lfloor i a_j/n\rfloor\). All the covers used here are connected,
as the branch exponents have greatest common divisor one with \(n\).
The rational idempotents of the cyclic deck group give Jacobian factors
up to isogeny. The factor for a divisor \(s>1\) of \(n\) has action by
\(\mathbb Q(\zeta_s)\), and its dimension is the sum of the coherent
character dimensions for characters of order \(s\). Complex conjugation
interchanges the coherent and holomorphic-one-form characters, preserving
the unordered multiplicities of each conjugate pair.
Ordinary reduction makes Frobenius invertible on coherent \(H^1\).
It suffices for the primitive factor to be ordinary because the
character-one line belongs to that factor. The nonzero scalar then
excludes the face by \eqref{eq:prime-isolation} and \eqref{eq:face-scalar}.
Work away from the finitely many primes of bad reduction for the curve,
cover, and deck action. Adjoin the required roots of unity and take
ordinary primes of absolute degree one, so they lie in the required
prime progression.

For heights three and four the compressed angular sets, up to sign, are
subsets of \(\{-2,-1,0,1\}\) or \(\{-1,0,1\}\). The branch exponents
are displayed below. The first two cases of the cubic partition are
all distinct and one double root. The third is a triple root.
\[
\begin{array}{c|c|c|c}
H & \deg P & \text{distinct roots} & \text{coincident roots}\\ \hline
3 & 2 &(2,1,1,2)&(2,2,2)\\
3 & 3 &(1,1,1,1,2)&(1,2,1,2),\ (1,3,2)\\
4 & 2 &(3,1,1,3)&(3,2,3)\\
4 & 3 &(2,1,1,1,3)&(2,2,1,3),\ (2,3,3)
\end{array}
\]
At height three the distinct quadratic and cubic root configurations
give genus two and three, respectively. In genus three the Jacobian has imaginary quadratic action
by \(\mathbb Q(\zeta_3)\). At height four the distinct-root quadratic
configuration has genus three, with an elliptic order-two factor and a
surface with order-four action. The distinct-root cubic configuration
has genus four, with an elliptic order-two factor and an imaginary
quadratic threefold. All repeated configurations have genus at most two,
except the triple cubic root at height three whose raw exponent equals
three. There \eqref{eq:face-scalar} is a unit times \(\binom{3k}{2k}\).
The genus at most two cases have ordinary primes of density one after
finite extension \cite{sawin}. For the genus-four case, first pass to a
common extension containing the roots of unity over which the elliptic
factor has density-one ordinarity. Apply Fit\'e's theorem over that same
field to the threefold with imaginary quadratic multiplication.
Its positive-density ordinary set intersects both the elliptic ordinary
set and the absolute degree-one primes in a set of positive density.
The same theorem gives the required primes for the genus-three case
\cite[Theorem~1]{fite}.

The noninteger heights with \(n\in\{7,9,10,11,15\}\) have quadratic \(P\).
For distinct roots the primitive signatures in increasing character order
are as follows.
\[
\begin{array}{c|c|c}
n&\text{branch exponents}&\text{primitive signature}\\ \hline
7 &(4,2,2,6)&(1,1,2,0,1,1)\\
9 &(8,2,2,6)&(1,1,2,0,1,1)\\
10&(5,3,3,9)&(1,2,0,1)\\
11&(7,3,3,9)&(1,1,2,0,0,2,2,0,1,1)\\
15&(10,4,4,12)&(1,1,0,2,0,2,1,1)
\end{array}
\]
In each row the primitive factor has rank two over \(\mathbb Q(\zeta_n)\)
and at most two conjugate pairs of multiplicities \((1,1)\), so
Theorem~\ref{lem:rank-two} applies. The deck action preserves the
Jacobian polarization, so its Rosati involution inverts the deck generator
and induces complex conjugation on \(\mathbb Q(\zeta_n)\). For \(n=9\)
there is also an elliptic factor of order three. For \(n=10\) the
additional factors have orders two and five, with dimensions one and
two, respectively. For \(n=15\) they have orders three and five, again
with dimensions one and two. There are no lower-order factors for
\(n=7,11\). If the roots coincide, merge the two equal interior
exponents. The character formula gives dimension zero at every lower
order in all five cases. The primitive factor has rank one over its
cyclotomic field. After
adjoining the field, its crystalline components at good degree-one primes
have height one and are ordinary. This proves exclusion in all five
rows, including the repeated-root cases.

For example, the height \(7/2\) face has angles \(-3,-1,1\).
Writing it as \(z^{-3}P(z^2)\), at \(p=7k+1\) the moment of order \(2k\)
detects \([z^{3k}]P(z)^{2k}\). The associated cover is
\(w^7=z^4P(z)^2\). When \(P\) has distinct roots, the Jacobian is a
sixfold with rank-two \(\mathbb Q(\zeta_7)\)-action. Its character pairs
\((1,6)\) and \((2,5)\) are balanced, while \((3,4)\) has multiplicities
\((2,0)\). The two balanced pairs are precisely the situation in which
the extra trace divisibility in Theorem~\ref{lem:rank-two} is used.

\subsection{The order-five collisions}

At height \(5/3\) a distinct-root quadratic gives exponents \((3,3,4)\) and
genus two. At height \(5/2\) the positive extreme is \(1\) and the negative
extreme is \(-3\) or \(-5\). In both cases distinct roots give exponents
\((4,2,2,2)\) and a rank-two order-five Jacobian, so
Theorem~\ref{lem:rank-two} applies. Merging two of the exponents \(2\)
gives genus two. There remain only the raw exponents exceeding five:
the repeated quadratic root at height \(5/3\) and the triple cubic root
at height \(5/2\).

For these cases use the sphere automorphism
\[
 \theta_{c,r}:(x,y,t)\longmapsto
 \bigl(x-cy^r t-c^2y^{2r+1}/2,\ y,\ t+cy^{r+1}\bigr).
\]
It preserves \(\Lambda\): its locally nilpotent derivation in the
Laurent \(y\)-chart is \(cy^{r+1}\partial_t\), and integration gives
zero by the endpoint formula, since \(x=0\) at \(t=\pm1\). Thus it also
preserves purity and safety by transporting markers.

At height \(5/3\), the face is \(Bx+Ay^2t+Cy^5\), with
\(A^2+2BC=0\) in the repeated-root case. It is the leading part of
\(B\theta_{-A/B,2}(x)\) for weights \((5,1,3)\). The inverse shear gives
\(Bx+P_4(y)+(a+b y)t\), where \(\deg P_4\le4\). The possible \(y^4\) term
is excluded by the two-node test. The next possible tie is between \(y^3\)
and \(yt\), since both have slope \(3/2\) from the node of \(Bx\). If both
occur, a distinct-root face gives the order-three cover with exponents
\((2,2,2)\), and a repeated-root face is the leading part of
\(B\theta_{-b/B,1}(x)\) for weights \((3,1,2)\). Invert this shear as well.
Any pure phase left has degree at most two.

At height \(5/2\), the triple-root face is the leading part of
\(B\theta_{c,1}(xt)\) for weights \((3,1,2)\). Indeed its face cubic is
\(B(z+c/4)^3\), as follows by expanding the shear and replacing
\(xy=(1-t^2)/2\). The inverse shear gives
\(x(Bt+a)+P_4(y)+tP_2(y)+d t^2\) with \(\deg P_i\le i\). The slopes from
\(Bxt\) to \(y^4,y^2t,y^3\) are \(7/5,4/3,5/4\), respectively, each strictly
largest once the earlier terms have been removed. The two-node test
removes all three. Again the remaining degree is at most two.

Section~\ref{sec:quadratic-height} gives safety for every pure phase of
degree at most two. This completes the unstable-leading branch and proves
safety for every pure phase of least sphere degree five.

\section{Proof of the ordinary-reduction theorem}\label{sec:ordinarity}

Fix a complex embedding of the base field. The rational Betti \(H^1\) is
two-dimensional over \(E\), and Hodge conjugation exchanges the embeddings
in each pair. The two holomorphic multiplicities of such a pair therefore
sum to two. We call the pair \emph{balanced} when both are one.

\begin{proof}[Proof of Theorem~\ref{lem:rank-two}]
Work over a finite extension containing a normal closure of \(E\), with the
endomorphisms and polarization defined. Choose a rational prime \(\ell\)
away from the discriminant of \(E\) and the index of the endomorphism
order in \(\mathcal O_E\). Take \(N\) with \(\ell^N>2^{2e+1}\) and further
extend the base to make \(J[\ell^N]\) rational and to contain the
\(\ell^N\)-th roots of unity. The absolute degree-one primes have density one,
since primes of higher residue degree have norm at least the square of
their underlying rational prime.
At a good such prime of residue characteristic \(p\), outside a finite
set, \(p\ne\ell\) and \(E\otimes\mathbb Q_p\) splits. The actual
endomorphism order in \(E\) has completion \(\mathcal O_E\otimes\mathbb Z_p\)
away from finitely many primes. Its idempotents split the \(p\)-divisible
group, and the crystalline summands have rank two.

Here are the trace comparisons used below. For an actual endomorphism
\(h\) of the reduction, the characteristic polynomial on \(H^1\) is
the degree polynomial of \(h\), both in \(\ell\)-adic and in crystalline
cohomology. Its value at a sufficiently large integer \(m\) equals
\(\deg([m]-h)\), by the exterior algebra description of cohomology
and the action on top cohomology. Take a rational basis of \(E\) made
of integral endomorphisms \(a\) and apply this to \(a\pi_p\), where
\(\pi_p\) is the \(p\)-power Frobenius of the reduction. The
nondegenerate field-trace pairing identifies the relative traces in
the two cohomologies as the scalar extensions of one element
\(\tau\in E\), determined by
\(\Tr_{E/\mathbb Q}(a\tau)=\Tr(a\pi_p\mid H^1)\).
The same argument with \(a\) alone identifies the ranks of the
crystalline summands with the rank-two Betti comparison ranks.
Every embedding of \(\tau\) gives a sum of two Frobenius eigenvalues,
which are algebraic integers of absolute value \(\sqrt p\) under
every complex embedding. Hence \(\tau\in\mathcal O_E\) and
\[
                             |N_{E/\mathbb Q}(\tau)|\leq 2^{2e}p^e.
\]

The split crystalline components have slopes in \([0,1]\), with sum
given by the dimension of the corresponding \(p\)-divisible factor.
The dimensions agree with the characteristic-zero Hodge
multiplicities in the split idempotent pieces, since the ranks in
the Lie algebra are preserved by good reduction. In each conjugate pair
of type \((0,2)\), the component of dimension two has both slopes one.
Its trace is divisible by \(p\).
In a nonordinary pair of type \((1,1)\), the two slopes of the
nonordinary component lie strictly between zero and one, since they
sum to one. The polarization pairs that component with slopes obtained
by subtracting its slopes from one. Both component traces are therefore
divisible by \(p\). The slopes show that each trace has positive
\(p\)-adic valuation, and integrality of the relative trace at each
split place makes that valuation at least one.
If \(b\) is the number of pairs of type \((1,1)\), nonordinarity implies
\[
                         p^{e-b+2}\mid N_{E/\mathbb Q}(\tau).
\]
Indeed all \(e-b+2\) places counted here are distinct and have degree one.
For \(b\leq2\) the normalized trace norm belongs to a fixed finite set
of integers at every nonordinary degree-one prime.

We now use a torsion congruence, as in the ordinary-reduction argument
attributed to Katz by Ogus \cite[Proposition~2.7.2]{ogus}. At our good
degree-one prime, Frobenius acts as the identity modulo \(\ell^N\) on
the Tate module, and \(p\equiv1\pmod{\ell^N}\). For every element \(a\)
of the integral endomorphism order in \(E\) we obtain
\[
  \Tr_{E/\mathbb Q}(a\tau)
    =\Tr(a\pi_p\mid H^1)
    \equiv\Tr(a\mid H^1)
    =\Tr_{E/\mathbb Q}(2a)
    \pmod{\ell^N}.
\]
The order completes at \(\ell\) to \(\mathcal O_E\otimes\mathbb Z_\ell\),
where the field-trace pairing is perfect by the choice of \(\ell\).
Thus \(\tau\equiv2\) in this completion modulo \(\ell^N\). Taking the
norm and using the congruence for \(p\) gives
\[
 p^{-e}N_{E/\mathbb Q}(\tau) \equiv 2^{2e}\pmod{\ell^N}.
\]
If reduction were nonordinary, the integer on the left would lie in
\([-2^{2e},2^{2e}]\). The modulus forces it to equal \(2^{2e}\).
Equality in the product Weil bound forces \(|\sigma(\tau)|=2\sqrt p\)
for every complex embedding \(\sigma\) of \(E\), so
\(\tau\overline{\tau}=4p\) in \(E\). At either split place of a
nonordinary balanced pair, the left side has valuation at least two
because both traces are divisible by \(p\). The right side has
valuation one for \(p\ne2\). This contradiction excludes all
nonordinary absolute degree-one primes outside the finite exclusions
over the chosen extension, proving the density claim.
\end{proof}

\subsection{Comparison with ordinary-reduction theorems}

Sawin determines the density of ordinary primes for abelian surfaces and
uses a Weil bound and algebraic Chebotarev to control an integral,
normalized Frobenius trace \cite[Theorem~1 and Corollary~2]{sawin}.
Fit\'e's results apply to certain threefolds, fourfolds, and higher-dimensional
varieties with extra endomorphisms. In particular, a Galois endomorphism
subfield of degree \(g=\dim J\) gives positive density of \(p\)-rank
at least \(\lceil g/2\rceil\) \cite[Theorems~1--4 and Corollary~5]{fite}.
These results supply several small-dimensional cases of the cover
argument.

Cantoral Farf\'an, Li, Mantovan, Pries, and Tang use a relative trace norm
to prove density-one ordinarity after finite extension for simple
signature with full geometric endomorphism algebra a CM field
\cite[Assumptions~2.1 and~2.5, Proposition~3.12, and Theorem~4.18]{cantoral}.
In relative rank two, simple signature permits one balanced conjugate
pair. Theorem~\ref{lem:rank-two} extends the potentially density-one part
of that rank-two result to two balanced pairs and to an acting field
\(E\) inside an arbitrary full geometric endomorphism algebra. The
additional divisibility is the trace contribution of the second
member of the nonordinary balanced pair. The proof works at split
primes after base extension. The results of \cite[Theorems~4.18--4.19]{cantoral}
also address higher relative rank, ordinary density over the original
field in certain cases, and \(\mu\)-ordinary reduction.

The primitive factors of cyclic orders \(7,9,11,15\) in the quintic
argument have dimensions \(6,6,10,8\), respectively, and the indicated
cyclotomic action has exactly two balanced pairs in each case. The
theorem applies to each smooth member with those branch exponents,
including members with additional endomorphisms. By comparison, the
order-five four-branch case has one balanced pair. Its exponents
\((4,2,2,2)\) become \((1,3,3,3)\) on changing the deck generator,
giving the Moonen family treated in
\cite[Corollary~5.1 and \S5.2]{cantoral} under the endomorphism hypotheses
there.

\section{A proper-volume obstruction}\label{sec:proper-volume}

The obstruction uses the fact that the original volume form extends
regularly to a proper model over the line of phase values. The
\(V\)-filtration records this regularity at the zero fiber. A nonzero
polynomial period would force the same volume class to have a nonzero
image in a module supported at zero, where its filtration degree is
impossible.

\begin{lemma}\label{lem:proper-volume}
Let \(Y\) be a smooth irreducible complex algebraic variety of dimension
\(n\), let \(f:Y\to\mathbb A^1_s\) be a proper nonconstant morphism, and
let \(\omega\) be a regular algebraic \(n\)-form on \(Y\). For any compact
closed \(n\)-cycle \(\gamma\) on \(Y(\mathbb C)\), the entire function
\[
 P_\gamma(\tau)=\int_\gamma e^{-\tau f}\omega
\]
cannot be a nonzero polynomial.
\end{lemma}

\begin{proof}
Factor \(f=\pi\circ i_f\), with \(i_f:Y\hookrightarrow
Y\times\mathbb A^1_s\) its graph and \(\pi\) the projection. Put
\(\mathcal B_f=i_{f+}\mathcal O_Y\) and
\(\mathcal M=\mathcal H^0(\pi_+\mathcal B_f)\), using left
\(\mathcal D\)-modules. Use the decreasing Kashiwara--Malgrange
filtration along \(s=0\), with
\(\partial_s s-\alpha\) nilpotent on \(\operatorname{gr}_V^\alpha\).
The graph generator \(e_f=\delta(s-f)\) belongs to
\(V^{>0}\mathcal B_f\). This is the graph-generator consequence of
strict negativity of the Bernstein--Sato roots
\cite[Corollary~5.2]{kashiwara-b-functions}
\cite[Corollary~5.15 and Example~5.19]{chen-dirks-mustata}.
In particular at each point of the graph over zero the generator
lies in \(V^\alpha\) for some rational \(\alpha>0\) locally.
Quasi-compactness of the fiber gives a common positive index
by taking the minimum over finitely many neighborhoods.
Off \(s=0\) every filtration piece is the whole graph module, and
off the graph there is nothing to check. Thus the top relative
form \(\omega e_f\) represents a degree-zero class in the relative
de Rham complex with coefficients in \(V^\alpha\mathcal B_f\).
The direct-image theorem for a smooth projection
\cite[Theorem~8.1]{chen-dirks-mustata} requires properness on the support
of the module. Here the support is the graph and the restriction of
\(\pi\) to it identifies with the proper map \(f\). The theorem therefore
identifies the image of the filtered complex with the corresponding
filtration piece of \(\mathcal M\). Consequently the class
\(m=[\omega e_f]\) satisfies
\[
 m\in V^\alpha\mathcal M\quad\text{at }s=0,\qquad \alpha>0.
\]
This step applies to the actual volume class even if the zero fiber
has a positive-dimensional critical locus.

We next check that the period respects the relations on the Fourier
transform of \(m\). Write \(M=\Gamma(\mathbb A^1_s,\mathcal M)\).
The graph action \cite[Equation~(27)]{chen-dirks-mustata} identifies
the graph module along \(Y\) with \(\mathcal O_Y[\partial_s]e_f\),
with actions
\[
 \begin{split}
 v(a\partial_s^j e_f)&=v(a)\partial_s^j e_f
                     -a\,v(f)\partial_s^{j+1}e_f,\\
 s(a\partial_s^j e_f)&=fa\partial_s^j e_f
                     -j a\partial_s^{j-1}e_f
 \end{split}
\]
for any local vector field \(v\) on \(Y\).
The Fourier convention is \(\tau=\partial_s\) and
\(\partial_\tau=-s\), which preserves the Weyl relation. Therefore
\(\widehat M\) is represented by degree-\(n\) hypercohomology on \(Y\) of
\[
 \bigl(\Omega_Y^\bullet[\tau],\,d_Y-\tau\,df\wedge\bigr),
 \qquad \partial_\tau\text{ acting by }\partial/\partial\tau-f.
\]
Indeed the relative de Rham complex of the projection uses the shift
by \(n\), and global sections on the affine target are exact.
One can compute this hypercohomology on a finite affine open cover of
\(Y\), whose finite intersections are affine since \(Y\) is separated.
This description uses the full polynomial parameter ring
\(\mathbb C[\tau]\).

At any \(\tau_0\in\mathbb C\), substituting \(\tau=\tau_0\)
in the total \v{C}ech complex commutes with both differentials.
After analytification, multiplying by \(e^{-\tau_0 f}\) turns the
specialized twisted differential into the ordinary holomorphic de
Rham differential. The affine intersections analytify to Stein
spaces, so the resulting \v{C}ech complex computes holomorphic de
Rham hypercohomology. Its class pairs with \(\gamma\) by analytic
de Rham comparison. A \v{C}ech coboundary thus has zero period
even when there is no global algebraic primitive. For a Weyl
operator \(A\), the class \(A\widehat m\) is represented by a global
polynomial family of top forms. Differentiating its exponential
integral on the fixed compact cycle realizes \(\partial_\tau\)
as \(\partial/\partial\tau-f\) on such a family. It follows that
evaluation on \(\gamma\) sends \(A\widehat m\) to \(A P_\gamma\)
and respects all relations on the cyclic submodule
\(\widehat N=\mathbb C[\tau]\langle\partial_\tau\rangle\,\widehat m\).

If \(P_\gamma\) were a nonzero polynomial, it would generate the
polynomial module \(\mathbb C[\tau]\) under Weyl operators, since
differentiation first produces a nonzero constant. The period
evaluation would give a surjection
\(\widehat N\to\mathbb C[\tau]\) taking \(\widehat m\) to \(P_\gamma\).
Inverse Fourier transform sends the polynomial module
\(D_\tau/D_\tau\partial_\tau\) to the point module
\(\delta_0=D_s/D_s s\), where \(D_u=\mathbb C[u]\langle\partial_u\rangle\).
It would give a quotient of a submodule of \(M\) sending \(m\)
nontrivially to \(\delta_0\). The induced filtration on submodules
and strictness for morphisms
\cite[Proposition~3.19(ii) and Corollary~3.20]{chen-dirks-mustata}
put the image of \(m\) in \(V^\alpha\delta_0\). For \(\alpha>0\)
the point module has zero such piece
\cite[Example~3.18]{chen-dirks-mustata}.
The contradiction proves the lemma.
\end{proof}


\paragraph{AI assistance and review.}
ChatGPT (OpenAI) assisted with mathematical arguments, literature checks,
exposition, proof review and verification code.

\begin{thebibliography}{99}

\bibitem{image}
A. van den Essen, D. Wright, and W. Zhao, On the Image Conjecture,
\emph{J. Algebra} \textbf{340} (2011), 211--224,
arXiv:\href{https://arxiv.org/abs/1008.3962v2}{1008.3962v2}.

\bibitem{cubicsphere}
R. van Rijn, \emph{Cubic sphere moments and the two-pair Image problem},
preprint, 2026.
\href{https://github.com/royvanrijn/jacobian-research/blob/57ca19131af2e3e3c6cea45abae9a94176d610f7/unpublished/cubic-sphere-image/cubic-sphere-image.pdf}{Versioned PDF}.

\bibitem{longcompact}
C. D. Long, \emph{The Mathieu Property for Compact Connected Lie
Groups}, arXiv:\href{https://arxiv.org/abs/2609.16178v1}{2609.16178v1}, 2026,
Theorem~3.1.

\bibitem{chen-dirks-mustata}
Q. Chen, B. Dirks, and M. Musta\c{t}\u{a}, \emph{An introduction to
\(V\)-filtrations}, arXiv:\href{https://arxiv.org/abs/2402.07816v1}{2402.07816v1},
2024.

\bibitem{kashiwara-b-functions}
M. Kashiwara, B-functions and holonomic systems. Rationality of roots of
B-functions, \emph{Invent. Math.} \textbf{38} (1976), 33--53.

\bibitem{ogus}
A. Ogus, Hodge cycles and crystalline cohomology, in
\emph{Hodge Cycles, Motives, and Shimura Varieties},
Lecture Notes in Math. \textbf{900}, Springer, 1982, 357--414,
\href{https://publications.ias.edu/sites/default/files/Number45.pdf}{Proposition~2.7.2}.

\bibitem{sawin}
W. F. Sawin, Ordinary primes for abelian surfaces,
\emph{C. R. Math. Acad. Sci. Paris} \textbf{354} (2016), 566--568,
arXiv:\href{https://arxiv.org/abs/1506.04784v2}{1506.04784v2}.

\bibitem{fite}
F. Fit\'e, Ordinary primes for some varieties with extra endomorphisms,
\emph{Publ. Mat.} \textbf{68} (2024), 27--40,
arXiv:\href{https://arxiv.org/abs/2005.10185v3}{2005.10185v3}.

\bibitem{cantoral}
V. Cantoral Farf\'an, W. Li, E. Mantovan, R. Pries, and Y. Tang,
\emph{Positive density of primes of ordinary reduction for abelian
varieties of simple signature}, arXiv:\href{https://arxiv.org/abs/2508.11174v1}{2508.11174v1},
2025.

\bibitem{muger-tuset}
M. M\"uger and L. Tuset, On the moments of a polynomial in one variable,
arXiv:\href{https://arxiv.org/abs/1910.12148v1}{1910.12148v1}, Theorem~2.1.

\end{thebibliography}
\end{document}
